% =====================================================================
% 01-Curvature.tex
% Geometric Hallucination Taxonomy — Mathematical Spine
% Borealis S. Hedling — 2026
% =====================================================================

\section*{01. Curvature Operator}

This module formalizes the curvature tensor, sectional curvature, and
instability criteria used to define Curvature–Driven Hallucinations (CDH)
in the Geometric Hallucination Taxonomy (GHT). All definitions rely on
the geometric foundations established in 00-Foundations.tex.

% ---------------------------------------------------------------------
\subsection*{1.1 Curvature Tensor}

Let $\nabla$ be the Levi–Civita connection associated with the metric $g$.
The curvature tensor is defined by
\[
R(X,Y)Z
= \nabla_X \nabla_Y Z
- \nabla_Y \nabla_X Z
- \nabla_{[X,Y]} Z.
\]

In local coordinates $(x^1,\dots,x^n)$, with Christoffel symbols
$\Gamma^i_{\;jk}$, the components of the curvature tensor are
\[
R^i_{\;jkl}
= \partial_k \Gamma^i_{\;jl}
- \partial_l \Gamma^i_{\;jk}
+ \Gamma^i_{\;mk} \Gamma^m_{\;jl}
- \Gamma^i_{\;ml} \Gamma^m_{\;jk}.
\]

The curvature tensor is multilinear and satisfies the usual symmetries:
\[
R(X,Y)Z = -R(Y,X)Z,
\]
\[
\langle R(X,Y)Z, W \rangle_g
= -\langle R(X,Y)W, Z \rangle_g,
\]
\[
R(X,Y)Z + R(Y,Z)X + R(Z,X)Y = 0.
\]

% ---------------------------------------------------------------------
\subsection*{1.2 Sectional Curvature}

Let $u,v \in T_p M$ be linearly independent tangent vectors. The sectional
curvature of the 2–plane $\Pi = \operatorname{span}\{u,v\}$ is
\[
K(u,v)
= \frac{\langle R(u,v)v, u \rangle_g}
       {\langle u,u \rangle_g \langle v,v \rangle_g
        - \langle u,v \rangle_g^2}.
\]

The denominator is the squared area of the parallelogram spanned by
$u$ and $v$, ensuring $K(u,v)$ is well–defined.

Sectional curvature measures how geodesics deviate from Euclidean
behavior in the direction of $\Pi$.

% ---------------------------------------------------------------------
\subsection*{1.3 Curvature Along Trajectories}

Given a latent trajectory $x(t)$, we define the curvature along the
trajectory by evaluating sectional curvature on the velocity vector
$\dot{x}(t)$ and a transverse direction $v(t)$:
\[
K_t = K(\dot{x}(t), v(t)).
\]

The choice of $v(t)$ depends on the diagnostic context. In synthetic
experiments, $v(t)$ is chosen to be orthogonal to $\dot{x}(t)$.

% ---------------------------------------------------------------------
\subsection*{1.4 Curvature Instability}

A curvature instability occurs when sectional curvature exceeds a fixed
threshold:
\[
|K(u,v)| > \kappa_{\mathrm{thr}}.
\]

Oscillatory instability occurs when the sign of $K(u,v)$ alternates
across steps:
\[
K_{t+1} \cdot K_t < 0.
\]

Collapse instability occurs when curvature approaches negative infinity:
\[
K(u,v) \to -\infty,
\]
indicating hyperbolic blowout.

Expansion instability occurs when curvature approaches positive infinity:
\[
K(u,v) \to +\infty,
\]
indicating spherical blowout.

% ---------------------------------------------------------------------
\subsection*{1.5 Curvature–Driven Hallucinations (CDH)}

Curvature–Driven Hallucinations arise when curvature instability causes
latent trajectories to diverge, collapse, or oscillate unpredictably.

Formally, a trajectory exhibits CDH if there exists a sequence of
directions $v(t)$ such that
\[
|K(\dot{x}(t), v(t))| > \kappa_{\mathrm{thr}}
\quad \text{for infinitely many } t.
\]

Equivalently, CDH occurs when curvature instability prevents the latent
trajectory from remaining within a stable geometric neighborhood.

% ---------------------------------------------------------------------
\subsection*{1.6 Curvature Stability Region}

Define the curvature stability region
\[
\Omega_{\kappa}
= \{ p \in M : |K(u,v)| \le \kappa_{\mathrm{thr}}
   \text{ for all } u,v \in T_p M \}.
\]

A trajectory is curvature–stable if $x(t) \in \Omega_{\kappa}$ for all $t$.

% ---------------------------------------------------------------------
\subsection*{1.7 Discrete Approximation}

In numerical settings, curvature is approximated using discrete
differences of Christoffel symbols:
\[
\partial_k \Gamma^i_{\;jl}
\approx
\frac{\Gamma^i_{\;jl}(x_{t+1}) - \Gamma^i_{\;jl}(x_t)}
     {\|x_{t+1} - x_t\|}.
\]

This approximation is used in the Python toolkit and validated in
\texttt{/tests/test\_curvature.py}.

% =====================================================================
% End of 01-Curvature.tex
% =====================================================================
